% =============================================================================
% Cover Letter — Quantum Science and Technology (IOP Publishing)
% Title: Quantum-Enhanced Agrivoltaic Digital Twin: Non-Markovian Coherence Control,
%        NPoM Plasmonic Diagnostics, and BB84-Secured IoT Telemetry
% Date: July 29, 2026
% =============================================================================

\documentclass[12pt]{letter}
\usepackage[utf8]{inputenc}
\usepackage[T1]{fontenc}
\usepackage{hyperref}
\usepackage[margin=2.0cm]{geometry}
\usepackage{siunitx}
\usepackage[version=4]{mhchem}
\usepackage{enumitem}
\DeclareSIUnit{\yr}{yr}
\DeclareSIUnit{\USD}{USD}
\DeclareSIUnit{\tonne}{t}
\DeclareSIUnit{\ha}{ha}
\DeclareSIUnit{\MWh}{MWh}
\DeclareSIUnit{\cubed}{^{3}}

\address{
Steve Cabrel Teguia Kouam\\
Department of Physics, Faculty of Science\\
University of Douala, Cameroon\\
\texttt{steve.teguia@facsciences-uy1.cm}
}

\begin{document}
\begin{letter}{Editor-in-Chief\\
\textit{Quantum Science and Technology}\\
IOP Publishing\\
Temple Circus, Temple Way\\
Bristol BS1 6HG, UK}

\opening{Dear Editor-in-Chief,}

We are pleased to submit our manuscript entitled \textbf{``Quantum-Enhanced Agrivoltaic Digital Twin: Non-Markovian Coherence Control, NPoM Plasmonic Diagnostics, and BB84-Secured IoT Telemetry''} for consideration as a Research Article in \textit{Quantum Science and Technology}.

\noindent\textbf{Publication Route}

\noindent We explicitly confirm that we opt to publish this manuscript via the \textbf{standard subscription route (carrying zero Author Processing Charge, \$0 APC)}.

\vspace{0.5em}
\noindent\textbf{Summary of the Work}

\noindent Agrivoltaic systems co-locating photovoltaics and crops conventionally rely on passive shading, which degrades agricultural yield. Here, we present a unified "quantum-to-organism" digital twin where an organic photovoltaic (OPV) layer integrated with a nanoparticle-on-mirror (NPoM) plasmonic cavity actively engineers the light spectrum to protect quantum coherence in the crop's Fenna--Matthews--Olson (FMO) light-harvesting complex, while simultaneously enabling single-molecule SERS diagnostics and quantum-secured telemetry.

\noindent Our main scientific contributions are:
\begin{itemize}[nosep]
\item \textbf{Non-Markovian Coherence Control \& Diagnostics:} Dual-band OPV filtering extends FMO excitonic coherence lifetimes by 50\% in the passive fraction and transport yields by 25\% at room temperature. We resolve the NPoM plasmonic transport trade-off, establishing an optimal mode volume ($V_{\rm mode} = 1.2\text{ nm}^3$) that preserves a global photosynthetic yield of 97.1\% across the canopy while enabling single-molecule SERS diagnostics.
\item \textbf{Greenhouse Microclimate \& Techno-Economics:} A 500~m$^2$ Cameroonian cooperative greenhouse model demonstrates a 28\% water saving (650~m$^3$/yr), 90~MWh/yr of electricity generation, and a 4.23-yr payback period under 30\% blended-finance grant support.
\item \textbf{Agro-Rural Cybersecurity:} Field telemetry is secured via room-temperature decoy-state BB84 QKD with $N_{\rm raw} \approx 1200$ bits, supplying low-latency (3.7~ms) edge inference deployed on embedded ARM Cortex-M4 IoT nodes.
\end{itemize}

\vspace{0.5em}
\noindent\textbf{Why \textit{Quantum Science and Technology}?}\nopagebreak

\noindent\textit{Quantum Science and Technology} is the premier venue for cross-disciplinary advances spanning quantum photonics, quantum biology, quantum sensing, and quantum communication. Our work bridges angstrom-scale non-Markovian vibronic dynamics (PT-HOPS/SBD) with field-scale agrivoltaic microclimates, demonstrating how quantum coherence control and quantum cryptography can deliver real-world sustainability and security benefits.

\vspace{0.5em}
\noindent\textbf{Prior Discussions and Related Submissions}

\noindent A companion manuscript addressing selective vibronic excitation in isolated FMO complexes is under revision at the \textit{Journal of Physical Chemistry Letters} (Manuscript ID: jz-2026-00994t). The present manuscript extends that quantum framework to a complete agrivoltaic digital twin incorporating NPoM cavity electrodynamics, SERS diagnostics, microclimate physics, lifecycle assessment, and BB84 cybersecurity. No part of this manuscript is under consideration elsewhere.

\noindent All authors have approved the manuscript and agree with its submission to \textit{Quantum Science and Technology}. The authors declare no competing financial or non-financial interests.

\vspace{0.5em}
\noindent\textbf{Suggested Reviewers:}
\begin{enumerate}[nosep]
\item \textbf{Prof. Jeremy J. Baumberg} --- Nanophotonics Centre, Cavendish Laboratory, University of Cambridge, UK (\texttt{jjb12@cam.ac.uk})
\item \textbf{Prof. Niek van Hulst} --- ICFO -- Institut de Ci\`{e}ncies Fot\`{o}niques, Barcelona Institute of Science and Technology, Spain (\texttt{niek.vanhulst@icfo.eu})
\item \textbf{Prof. Gregory D. Scholes} --- Department of Chemistry, Princeton University, USA (\texttt{gscholes@princeton.edu})
\item \textbf{Prof. Alexandra Olaya-Castro} --- Department of Physics and Astronomy, University College London, UK (\texttt{a.olaya@ucl.ac.uk})
\item \textbf{Prof. Joel Yuen-Zhou} --- Department of Chemistry and Biochemistry, University of California San Diego, USA (\texttt{joelyuen@ucsd.edu})
\end{enumerate}

\closing{Sincerely,}

\vspace{0.5em}
Steve Cabrel Teguia Kouam\\
(On behalf of all authors)

\vspace{0.5em}
\noindent\textit{Enclosures:} Manuscript (\texttt{Manuscript\_QST\_26-07-29.pdf}), Supporting Information (\texttt{SI.pdf}), and Figures.

\end{letter}
\end{document}
